\chapter{INTRODUCTION}
\pagestyle{fancy}
\pagenumbering{arabic}

\section{Research Motivation}
In modern data-driven organizations, streaming data pipelines process millions of events per second. Ensuring data quality in these pipelines is critical for downstream analytics, machine learning models, and business decisions. Unlike batch data processing where quality checks can be applied retrospectively, streaming environments demand real-time validation with minimal latency overhead.

The NYC Yellow Taxi dataset from the NYC Taxi and Limousine Commission (TLC) serves as our case study. Analyzing records from this dataset reveals five categories of data quality problems that existing streaming DQ frameworks cannot adequately address: (1) completeness violations, where required fields are missing or null; (2) validity violations, where individual field values violate physical constraints such as negative fares, zero distances, or impossible speeds; (3) multivariate anomalies, where combinations of individually reasonable feature values are collectively suspicious and require understanding inter-feature correlations; (4) uniqueness violations, where duplicate records appear due to data reprocessing; and (5) concept drift, where the underlying data distribution changes over time, causing previously calibrated thresholds and models to become stale. These problems demand real-time detection with minimal latency overhead, as streaming environments process millions of events per second without the luxury of retrospective batch analysis.

This motivates the CA-DQStream framework: a four-stage streaming pipeline that addresses each problem category with specific techniques, and monitors its own performance to detect when baselines become stale.

\section{Gaps in Current Research}
Despite decades of research in data quality management, three critical gaps remain in existing streaming DQ frameworks:

\textbf{Gap 1: Context Collapse in Heterogeneous Data.} Traditional frameworks apply a single global threshold to all records regardless of context, producing catastrophically high false positive rates on heterogeneous data.

\textbf{Gap 2: Sequential Pipeline Architectures Limit Performance.} Existing frameworks process records through validation stages sequentially, compounding end-to-end latency and wasting resources on records that earlier stages could flag.

\textbf{Gap 3: Single-Strategy Drift Adaptation.} When concept drift is detected, existing frameworks apply full model retraining regardless of drift severity, wasting compute on minor shifts that lightweight updates could resolve.

These gaps are detailed in Chapter 3, where we derive five research questions and six hypotheses to address them.

\section{Expected System}
This thesis introduces \textbf{CA-DQStream} (Context-Aware Data Quality for Streaming), a framework organized into five layers, built on Apache Kafka \citep{kafkaintro2024} for distributed event streaming and Apache Flink \citep{flinkdocs2024} for stateful stream processing with checkpoint-based fault tolerance:

\begin{enumerate}
    \item \textbf{Data Source Layer} --- NYC TLC Parquet files ingested via Kafka Producer into the streaming pipeline.
    \item \textbf{Kafka Layer} --- Distributed event streaming backbone with six topics managing raw data, violation records, anomaly scores, meta-streams, and model updates.
    \item \textbf{Processing Layer with Flink} --- A four-stage Rendezvous pipeline implementing the scientific contributions of this thesis: Schema Validation $\rightarrow$ [Business Rules $\parallel$ Isolation Forest] $\rightarrow$ MetaAggregator $\rightarrow$ IEC (ADWIN Drift Detection). The fork-join pattern parallelizes rule-based and ML-based validation, enabling early exit optimization.
    \item \textbf{Data Storage Layer (MinIO)} --- Persistent storage for cleaned data, dirty data, and alerts in Parquet format, accessible for downstream analytics.
    \item \textbf{Monitoring Layer (Prometheus + Grafana)} --- Real-time observability into pipeline health, violation rates, anomaly scores, and drift alerts.
\end{enumerate}

The Monitoring Layer complements the processing pipeline by providing full observability: Prometheus scrapes metrics from Flink TaskManagers, MinIO, and Kafka; Grafana dashboards display DQ Overview (violation rates per layer), ML Comparison (anomaly scores by priority), System Performance (throughput, latency), and Drift Alerts (ADWIN events).

\section{Objectives and Contents}
The primary objectives of this thesis are:

\begin{enumerate}
    \item Design and implement a five-layer architecture for streaming DQ: Data Source, Kafka, Processing with Flink (Schema Validation, Business Rules, Isolation Forest Anomaly Detection, ADWIN Drift Detection), Data Storage, and Monitoring.
    \item Develop a four-stage processing pipeline: Schema Validation, Business Rules, Isolation Forest anomaly detection, and ADWIN drift detection.
    \item Develop a self-monitoring mechanism using ADWIN drift detection on quality meta-streams (MetaAggregator $\rightarrow$ dq-meta-stream $\rightarrow$ ADWIN).
    \item Evaluate CA-DQStream on the NYC Taxi dataset demonstrating improved false positive rate control and effective drift detection.
    \item Implement MLOps pipeline with Isolation Forest model training, MLflow tracking, and in-memory deployment via Flink Broadcast State.
\end{enumerate}

The thesis is organized as follows: Chapter 1 covers background and related works; Chapter 2 presents the problem definition and research questions; Chapter 3 details the four core innovations of CA-DQStream; Chapter 4 describes the system architecture and implementation; Chapter 5 presents experiments and evaluation results; and the final chapter concludes with contributions and future directions.

\section{Key Contributions}
The key contributions of this thesis are:

\begin{enumerate}
    \item \textbf{Five-Layer Architecture:} A production-grade streaming DQ architecture built on Apache Kafka and Apache Flink, with persistent storage via MinIO (S3-compatible object storage) and full observability through Prometheus and Grafana. This architecture separates infrastructure concerns from algorithmic contributions, enabling independent optimization of each layer.

    \item \textbf{Four-Stage Rendezvous Pipeline:} A novel fork-join processing pipeline that parallelizes rule-based (Canary) and ML-based (Complex) validation branches, converging at a MetaAggregator rendezvous point. This architecture reduces end-to-end latency and increases throughput compared to sequential pipelines.

    \item \textbf{Context-Aware Thresholding Framework:} A four-dimensional threshold matrix that assigns different anomaly detection thresholds to different data contexts, preventing false positives from legitimate high-value transactions that appear anomalous under global thresholds.

    \item \textbf{Intelligent Evolution Controller:} A multi-strategy self-adaptation mechanism that monitors the quality meta-stream via ADWIN and selects from four complementary adaptation strategies (Continuous Evolution, Switching, METER, Spatial Tracking), resolving most drift events through lightweight incremental updates rather than expensive full retraining.

    \item \textbf{Hybrid K-Means + sklearn IsolationForest Model:} The integration of K-Means clustering with sklearn IsolationForest anomaly detection, where clustering provides context boundaries and Isolation Forest scores anomalies within each context cluster. This architecture detects multivariate anomalies that single-feature business rules cannot express, such as meter tampering and short-trip fraud.
\end{enumerate}

\section{Thesis Organization}
The remainder of this thesis is organized as follows:
\begin{itemize}
    \item \textbf{Chapter 1: Introduction} --- Provides an overview of streaming data processing challenges, introduces the CA-DQStream framework, and outlines the thesis structure.
    \item \textbf{Chapter 2: Background and Related Works} --- Reviews streaming data processing (Kafka, Flink), MLOps, DAMA-DMBOK DQ dimensions, and related works on context-aware DQ, Rendezvous/fork-join patterns, and drift detection.
    \item \textbf{Chapter 3: Problem Definition and Research Questions} --- Presents empirical analysis of the NYC Yellow Taxi dataset as a case study, research gaps, five research questions (RQ1-RQ5), and six hypotheses (H1-H6).
    \item \textbf{Chapter 4: Core Innovations} --- Details the four scientific contributions: (1) 4D Context-Aware Thresholding, (2) Rendezvous Pipeline Architecture, (3) Intelligent Evolution Controller with 4 Adaptation Strategies, and (4) Hybrid K-Means + sklearn IsolationForest Model.
    \item \textbf{Chapter 5: System Architecture \& Implementation} --- Describes infrastructure (Kafka, Flink, MinIO, Grafana), four-layer processing pipeline, optimization techniques (PyFlink Vectorized UDFs, StreamingFileSink), and MLOps integration.
    \item \textbf{Chapter 6: Experiments \& Evaluation} --- Presents 13 experiments organized by research question: RQ1 (multi-layer coverage), RQ2 (Hybrid model vs baselines), RQ3 (IEC drift adaptation), RQ4 (Rendezvous performance), RQ5 (4D thresholds vs global), plus additional validation experiments.
    \item \textbf{Chapter 7: Conclusion \& Future Work} --- Summarizes contributions, validates hypotheses, discusses limitations, and outlines future research directions (multi-domain validation, deep learning integration, federated learning, anomaly explanation via intrinsic structure).
\end{itemize}
